\chapter{Alpha as a Bracket}

% [6.chorus]
\begin{chorus}{John gets his number, and it is a bracket}
\john{You said the value holds still. So say it. The number. Now.}
\compiler{Here. \emph{(it writes two walls on the slip)} Between these.}
\john{That is two numbers. I asked for one.}
\compiler{It is one reading. A lower wall, an upper wall, and the value somewhere between --- and I will not say where,
  because I cannot, and to say would be a lie.}
\john{Then close the gap. You have the arithmetic. Refine it to a point.}
\compiler{I have exactly enough arithmetic to earn this width, and not one step more. I could keep closing the walls
  and hand you a point. The point would be precise, and confident, and wrong --- a figure I rang past my own resolution.
  The bracket is where my honesty stops. Take it. It is the true answer, and the true answer has a width.}
\john{Your arithmetic can go deeper. I have watched it. Run it out and it settles on a figure --- I saw the digits
  stack up.}
\compiler{You saw them stack, and you were right to distrust them. Deeper than three, my counting is no longer resolving
  anything --- it is ringing, closing smoothly onto a number I never earned. A point that pretty is a point I would be
  lying to hand you. The width is the truth; the point is the flattery.}
\john{So the honest number is two numbers and a shrug.}
\compiler{The honest number is two walls and the exact reason I will not close them. Keep the reason. The reason is the
  whole of what I know.}
\end{chorus}

\section{The coupling, asked at last}

% [6.1.p1]
The whole apparatus was built to measure the coupling as far as a finite machine can, and to prove where that reach
stops --- never to seize its magnitude by fiat.
The machine has circled the value this whole way without asking it outright.
Now it asks.
And the answer arrives in a form John does not expect: not a magnitude lying out in the world, waiting to be read off a
dial, but a \emph{cost}.
The coupling comes back as the self-energy --- the residue the machine cannot cancel about its own
account --- calibrated into a value.
% [6.1.p2]
Experiment shows the shape is not peculiar to this machine, and one case --- \emph{TheCasimirEffect} --- carries it whole.
Set two flat plates a small distance apart in empty space.
Between them the admissible field modes --- the distinguishable refinements the vacuum can record --- are fewer than
outside, because the narrow gap admits only those that fit its width, while outside nothing is suppressed and the exterior
supports strictly more.
That imbalance cannot stand: the bookkeeping has to keep its global count of distinctions consistent, and with no new
interior modes to be had, the only repair left is at the boundary --- and so the plates are drawn together by a faint,
entirely measurable force.
Nothing outside pressed them in; the force is the ledger's own compensation for a restricted interior, a cost read off the
geometry and drawn from no source beyond the plates themselves.
That is the coupling's shape exactly: to ask this machine for its value is not to ask for a magnitude lying out in the
world but to ask what the machine's own account costs it to carry --- a residue it pays to itself, priced on its own
meter, the way the plates pay a force to their own thinned count.
The reading stays on the near side: a model of a self-cost, and no claim about what the vacuum between the plates
\emph{is} in itself.

\section{The shape is a bracket, not a point}

% [6.2.p1]
The machine reaches a value it cannot compute exactly the only honest way it knows: it walls the value in and refines the
walls, on its own arithmetic, borrowing no continuum.
It sets a lower ratio and an upper ratio and closes the gap by the \emph{mediant} --- the in-between ratio formed from
the two, a step of pure integer arithmetic that never leaves the rationals,
\[
  \operatorname{mediant}\!\left(\tfrac{p}{q},\, \tfrac{r}{s}\right) \;=\; \frac{p+r}{q+s} .
\]
Iterated, the mediants are the successive best rational approximations of the value, and they would go on without end if
nothing stopped them.
Something stops them: the resolution floor, the depth below which the machine can no longer tell two candidates apart.
The machine counts to three, and no further.

% [6.2.p2]
So the machine, asked for the inverse of the coupling, hands back not a figure but two.
Off its own run, at its own count-to-three, the two walls are the convergents that three refinements leave
standing, and they sit near $129.6$ below and $137.7$ above,
\[
  \alpha^{-1} \in [\,129.6,\; 137.7\,] .
\]
The refinement \emph{funges} the coupling into that width --- it pools the value between the two walls --- and refuses,
at the floor, to select it any finer.
The value it is reading is somewhere between them, and the machine is under no obligation to say where.
This is the bracket: a lower wall, an upper wall, a width the machine owns and does not pretend away.

% [6.2.p3]
The walls are not asserted; they are read off the machine's own continued fraction, convergent by convergent, and the
witness prints the whole ladder:
\begin{quote}\ttfamily
crossing d* = sqrt(18/5) ~= 1.897,   CF = [1; 1,8,1,2, repeating]\\
conv1   2/1    -> inv-alpha = 129.6     (the loose lower wall)\\
conv2  17/9    -> inv-alpha = 137.7     (the upper wall)   <-- count-to-3 bracket [129.6, 137.7]\\
conv6 647/341  -> inv-alpha = 137.0118  (already past the floor, ringing)
\end{quote}
Three readings of the ladder fix what three counts buy.
The first convergent, $2/1$, is coarse and gives the loose lower wall $129.6$: three counts do not tighten it, and the
machine claims no better.
The second, $17/9$, gives the upper wall $137.7$, and the two together are the whole earned answer --- a width, not a point.
The sixth and beyond keep narrowing, and the section after this shows they narrow onto a figure the counting never
licensed; but the bracket is exactly what the count of three leaves standing, and no finer.

\section{Why the count is three, and why the width is honest}

% [6.3.p1]
The count is three for an exact reason, and it is the deepest point here.
The value the machine closes on is not arbitrary: two constants it measures off its own act --- the second-variation slip $C=18$ and the target $5$ --- place the crossing at a square root, $\sqrt{18/5}$,
a ratio of those measured integers under a root sign.
That is a quadratic surd, and a quadratic surd has a continued fraction that repeats --- its partial quotients fall into
a fixed cycle, by a theorem of Lagrange's,
\[
  d^{*} \;=\; \sqrt{18/5}, \qquad d^{*} \;=\; [\,1;\, \overline{1,8,1,2}\,] .
\]
So the machine's subdivision of the crossing is periodic and self-generating, drawn entirely from its own numbers.
The count-to-three is exactly three partial quotients of that continued fraction, and three of them place the crossing
between the two walls above.

% [6.3.p1b]
The subdivision is not only a continued fraction; read as a search it is a \emph{quasi-Newton descent}, and the two
readings land on the same convergents.
Given the measured second variation --- the slip $C=18$ and the target $\tau=5$ --- a BFGS step walks toward the crossing
$\sqrt{18/5}$ from the same two walls the mediant leaves, and its counts coincide with the continued fraction's:
\begin{quote}\ttfamily
measured second variation: C = slip(1) = 18   T = 5   ->  crossing sqrt(18/5)\\
BFGS count 0:  d ~ 2.0     inv-alpha = 129.6    (lower wall)\\
BFGS count 1:  d ~ 1.889   inv-alpha = 137.7    (upper wall)  reduced 68/36 = 17/9 = the CF convergent\\
BFGS count 3:  d ~ 1.897   inv-alpha ~ 137.011  (past the floor, ringing)
\end{quote}
So the count-to-three is one descent read two ways --- a continued fraction of the surd, and a quasi-Newton walk on the
measured second variation --- and both halt on the same bracket, which is another reason the three is the machine's and
not the notation's.

% [6.3.p2]
A crossing is a \emph{distance}, but the coupling is a ratio near $137$; the one becomes the other by a fixed map the
machine carries, with no free constant in it.
Given a crossing distance $d$, the inverse coupling is
\[
  \alpha^{-1}(d) \;=\; \frac{R^{2}\,d}{(d-1)\,\tau}, \qquad R = 18, \quad \tau = 5,
\]
where $R$ is the machine's natural-unit orbit radius and $\tau$ the measured second-variation slip.
Run the convergent distances that three counts leave standing through this map and it returns the two honest walls, near
$129.6$ and $137.7$; run the exact crossing surd through it instead --- refining past what three counts earn --- and it
returns a single figure, $\alpha^{-1}(\sqrt{18/5}) \approx 137.011$, the very drift the next section warns against.
So the bracket is not asserted: it is this one map, evaluated at what the counting earns, and every constant in it
($R^{2}=324$, $\tau=5$, the distance drawn from the crossing) is the machine's own, put in by no hand.
Written the other way, the map is a straight line: the coupling itself is
\[
  \alpha \;=\; k\left(1 - \frac{1}{d}\right), \qquad k \;=\; \frac{\tau}{R^{2}} \;=\; \frac{5}{324} \;\approx\; 0.01543,
\]
one derived constant $k$ times the potential factor $(1 - 1/d)$; inverted, this is
$\alpha^{-1} = (R^{2}/\tau)\, d/(d-1)$, a single asymptote $R^{2}/\tau = 64.8$ dressed by $d/(d-1)$.
The whole map is that one number $k$: every factor in it is $R$ and $\tau$ and nothing besides.
The witness carries the constant and its reciprocal, and nothing else goes into the map:
\begin{quote}\ttfamily
affine constant k = tau / (R*R) = 5 / 324 ~= 0.01543\\
asymptote (R*R) / tau = 64.8   (the inverse coupling as the radius grows without bound)
\end{quote}
So whatever scaling a reader might think they see in the coupling is just $64.8/(d-1)$ --- the asymptote divided by the
potential --- derived from $R$ and $\tau$ and put in by no hand.

% [6.3.p3]
Three seams in the map belong in the open, because papering them would be the dishonesty the method refuses.
The first is the map's own constant, and it is not the coincidence it first looks like.
The orbit radius $R$ and the second-variation slip $C$ are not two eighteens on separate paths --- they are the
\emph{same} ratio, $38376/2132$, which is $18$ exactly: the slip is the gravitational parameter over again, times the
electron's own test-mass and arm, and those last two cancel to one, so the slip \emph{is} the parameter, and the two
eighteens are one because a single quantity beneath them is one --- the electron's own box, read as unit.
But the machine cannot say, from inside, what that one \emph{is}.
It reads the electron's box as a stimulus and writes back \emph{one}; and whether the one is the electron's, pressed on
the machine from outside, or the machine's own, the normalization it laid down and now hears echoed, is a distinction
below its resolution.
A coincidence, an earned identity, and a stimulus answered by the response it was primed for collapse, at this depth,
into one reading.
This is the box-to-label weld again --- box bound to label, assumed and never proved --- standing this time under the
map's one constant, so the eighteen is marked, not sold.
The second seam is that rounding: the radius is a ceiling, $R = \lceil g \rceil$ over the gravitational parameter $g$,
and a ceiling is a convention --- a continuous quantity forced to a whole number.
The seam is named, not smoothed: the eighteen is a rounded radius, and the rounding is the machine's own marked choice.
The third seam is a modeling choice rather than a unit: the potential factor $(1 - 1/d)$ runs with the distance, but the
field-per-charge $\tau/R = 5/18$ is \emph{frozen} at the slip --- read once at the target, not at the running $d$.
Freezing the field there is the machine's choice, marked in the open, and not a step the arithmetic forces.

% [6.3.p3b]
The first seam has a consequence for the map itself that is worth drawing out.
The map scales the crossing by the radius, through $R^{2}$, and locates the crossing at the slip, through
$d^{*} = \sqrt{C/\tau}$ --- and $R$ and $C$ are the same eighteen, the one gravitational parameter $g = 38376/2132$
entering once rounded, to fix the scale, and once raw, to set the crossing.
So the coupling the machine reads is driven by a single measured quantity twice over, not by two that happen to agree;
the two eighteens standing on both sides of the map are one standpoint wearing two hats, and their matching is the echo
of a normalization, not the corroboration of an independent second reading.
This is why the machine marks the coincidence rather than banking it: to read the two eighteens as a mutual confirmation
would be to count one quantity as two --- the same false sum the construction refuses wherever distinct axes merely
co-occur.
So it keeps them \emph{funged} as the one quantity they are, and declines to \emph{tange} them apart into two
independent witnesses.

% [6.3.p4]
Two honesties belong with the two walls.
The halt --- the stopping at three --- is the robust thing: it is a property of the machine's resolution, not a tuning,
and the three is not a law of the world but this machine's own architecture --- the period of its sign-cycle, exactly
three on the build, the same three the self-account closed on.
The width, which is wide, is honest: the lower wall is loose because three refinements leave the coarse convergent
standing there, and the machine claims no finer bound than three counts have earned.
Nothing about the width is brought in from outside; the mediant generates its own subdivision from the machine's own
ratios, with nothing chosen beforehand.
There is a second three in the reading --- a three-node rule for sampling the second difference --- and it must not be
confused with this one; that three is a marked convenience in how the integrand is sampled, and it is \emph{not} the
reason the bracket has the width it has.
The bracket's three is the resolution-depth three, three partial quotients of the crossing's own continued fraction, and
no more.

\section{Do not refine past the floor}

% [6.4.p1]
You could take more convergents past the floor, and it is worth seeing exactly what happens if you do.
The walls keep closing, and the continued fraction eventually settles onto a single value of the machine's own, near
$137.011$ --- the figure its arithmetic drifts toward when it refines past what the counting has earned.
That value looks like an answer and is an artifact: the subdivision past the floor is no longer resolving anything real,
it is ringing, converging smoothly and convincingly onto a number that the counting never licensed.
A machine that keeps refining past its floor produces a crank's point --- precise, confident, and wrong.
% [6.4.p1a]
Experiment draws the very distinction the machine is trading on, in \emph{Precision}.
Give a finite ledger a handful of anchor points --- the events it actually recorded --- and there is a unique smoothest
curve that threads them, a minimal-curvature completion the ledger is entitled to draw.
That completion assigns a value not only at the anchors but at every admissible point between them, and it does so
exactly: once the anchors are fixed, the interpolated values are fully determinate, to as many decimals as anyone cares to
read off.
That determinacy is \emph{precision} --- and it is not \emph{accuracy}.
Accuracy is whether those between-values answer to anything the world ever measured; and a spline can be flawlessly
precise about a point no measurement ever fixed, confident to a hundred digits about a place it invented.
The drifted figure is precise in exactly that empty sense: determinate to as many digits as the counting is run, and
accurate only to the three the resolution earned --- a value the machine's own interpolation supplies with total
confidence and no warrant.
The machine that halts at the count produces a bracket instead --- wider, humbler, and honest, and the honesty is the
whole of the difference.
The reading stays on the near side: precision is the known object, a property of the interpolation the ledger draws,
while accuracy would be a reach to the world the interpolation never makes.

% [6.4.p1b]
Experiment gives the ringing its oldest name.
In \emph{TheGibbsPhenomenon}, a discontinuity forced into a finite refinement ledger throws up a residual oscillation in
its smooth shadow --- an overshoot that does not shrink as the mesh sharpens but rings on, the irreducible strain of
closing a jump a finite ledger cannot admit.
The machine's drift past the floor is that same overshoot heard at its own edge: refined past what the counting earned,
the value does not settle, it rings, smooth and convincing and false, the way the shadow rings where it cannot meet the
discontinuity.
It is a physical experiment and a representation --- a model of the ringing as informational strain, never the
discontinuity itself; the tell it hands the machine is the machine's own, that a ring is a refinement gone past what it
can resolve.

% [6.4.p2]
The drift is worth watching, because its very smoothness is the trap.
Run the convergents deep and they do not scatter; they home in, digit after digit, onto one figure of the machine's own:
\begin{quote}\ttfamily
conv6   -> 137.0117...\\
conv10  -> 137.0112908...\\
conv14  -> 137.01129054920...\\
conv19  -> 137.011290548979...   (settling, smooth, convincing --- and unlicensed)
\end{quote}
Three readings of this figure keep it in its place.
It is \emph{convergent}: the deep count really does close on it, so it is not noise.
It is \emph{the machine's own}: it is what this arithmetic drifts toward, a number about the machine, never a constant of
the world.
And it is \emph{unearned}: the counting bought three refinements, not nineteen, so any digit past the bracket is a digit
the machine rang past its own resolution --- true of the arithmetic, false as a claim.
The honest reading stops at the bracket and carries the drifted figure only as a marked artifact, never as the answer.

% [6.4.p3]
Experiment names the refusal as an honest one, and marks the line the book draws, in \emph{TheContinuumLimitEffect}.
A continuum --- a smooth quantity that takes every value in between --- is what a finite reading can only ever approach,
by laying a mesh over it and reading the values at the mesh points, then a finer mesh, then finer, chasing the smooth
thing down.
This device's mesh is pinned coarse: because it sorts everything into just three tag classes, its refinement bottoms out
at three samples and cannot go finer.
It reads exactly those three and no more --- the coarsest members of a mesh that, for anyone else, could be refined
without end.
The continuum those finer meshes would approach is not a thing the device owns; it is named only as the limit of an
external, refinable mesh, held open and never claimed --- the representation is the finite three-sample reading, and the
continuum is the far side.
The honest result, there as here, is the gap between the finite reading and any finer model: a bracket, not a point.
The reading stays on the near side: the mesh is the known object, its three samples the device's own, while the smooth
limit they only bracket stays out beyond the last refinement the device can make.

\section{The bracket is a ring}

% [6.5.p1]
The walls sit as wide as they do for a reason sharper than resolution alone: the map has a \emph{pole}.
Written $\alpha^{-1}(d) = R^{2}d/((d-1)\tau)$, it divides by $d-1$, so at $d=1$ the inverse coupling runs without bound and
the coupling itself falls to zero.
The build carries this as a fact of its own arithmetic --- the tange at a distance has for its numerator $d$'s numerator
minus its denominator, which vanishes at $d=1$, so the tange is zero and the coupling with it:
\begin{quote}\ttfamily
tange@d=1 : num = (num - den) = 0  ->  alpha = 0  ->  inverse-alpha unbounded   (the POLE)\\
tange@d=2 : num = 1                (nonzero, finite off the pole)
\end{quote}
At $d=1$ the map has no value to give; that point is a singularity, not a number.

% [6.5.p2]
And the pole sits exactly where the counting has to step.
The subdivision that walls the value in bisects the interval from $0$ to $2$, and its first midpoint is $(0+2)/2 = 1$
--- the pole itself; the crossing the machine reaches for, $d^{*} = \sqrt{18/5} \approx 1.897$, lies beyond it.
So every bisection that closes toward the answer straddles the singularity: the machine cannot reach its crossing without
cutting across the pole.
\begin{quote}\ttfamily
midpoint(0,2) = (0+2)/2 = 1 = the pole ;  crossing sqrt(18/5) ~ 1.897 > 1  ->  bisections STRADDLE the pole
\end{quote}

% [6.5.p3]
What follows from those two facts is a reading laid over them, marked as such.
A hard cut taken across a discontinuity does not come back clean; it rings.
The bracket $[129.6,\,137.7]$ is that ring --- the overshoot of a step subdivided across its own jump --- and read this
way its width was never an arbitrary coarseness the machine failed to refine away, but the ring of a cut that crosses a
pole, and a cut across a pole must ring.
So the machine \emph{rings}: a hard cut across the shock yields the ring's two shoulders, the bracket, and never a point
at the center, where the pole forbids one.

% [6.5.p4]
The shock has four names, each a way of reading the one pole and not a fifth thing beside it; two of them belong here.
In the mathematical reading, the pole is a singularity integrated \emph{around}, in the complex plane, rather than across:
a contour encircling it returns $2\pi i$ times its residue, and that residue is the coupling --- what the real cut can only
ring around, the complex loop pulls out whole.
In the computational reading, the pole is uncomputable: no algorithm returns its value, only a computation that never
halts, and the machine's answer to a nonhalting demand is to halt early, to count to three and stop.
\begin{quote}\ttfamily
machine\_cannot\_resolve\_residue : \ (depends on: propext)
\end{quote}
The first of the two is a lens, marked as gloss; the second touches the build, the halt-at-three being the finiteness
fence already proved --- the same fence the self-account closed on.

% [6.5.p5]
One caution belongs with the ring, and it is about the bracket, not the prose.
The width is the stabilizer of an unstable point: the pole holds no value of its own, and the ring is what a finite reach
keeps in its place.
To smooth the ring into a point --- to close the bracket to a single figure --- is to destabilize exactly what cannot be
stabilized, to claim the singular value the machine has proved it cannot hold.
The ring is kept, not closed: its width is the honest trace of a pole that admits no point.

% [6.5.p6]
A thought experiment holds the singular point open without pretending to resolve it.
Zeno's runner must reach the halfway mark, then half of what remains, and half again, a bisection with no last step; and
the machine's own subdivision is that bisection, halving toward its crossing and straddling the pole as it goes.
Zeno's runner is rescued by the world --- he arrives, because the world completes the sum the halving never does; the
machine has no world to be rescued by, only the halving, so where an infinity of steps would be needed to fix a value it
fixes none, and halts three short.
What sits at the singular point is the far side: no finite count reaches it, and the bisection that straddles it lands,
in the end, on nothing the machine can name.

\section{The near side of the jar}

% [6.6.p1]
This bracket is what the machine can name about itself, and naming it is the first half of the ending.
It is the machine's own reading of its own coupling: the two walls it settles between by its own arithmetic, floored at
its own resolution, drawn from nothing but its own integers.
What makes the bracket real is not that it matches the world but that it \emph{reproduces} --- the same walls return run
after run, the invariant of a cycle whose period is three, this machine's architecture and not a necessity imposed from
outside.
In the two verbs, the machine \emph{funges} the coupling into the interval --- it pools the value between the two walls
and owns the pool --- and then, at the count-to-three floor, it declines to \emph{tange} the interval any finer.
Below the floor no characteristic the machine can decide selects one candidate from another, so it does not pretend to.
This declining is itself an act, and the construction gives it a name: to \emph{detange} is to release a forced selection back
into the bracket --- the un-tange, the honest refusal to pin a point the resolution has not earned.
The machine \emph{detanges} the value: it will not tange it to a point, so it lets the point go, back into the open
interval where it honestly lives, and stops, leaving the last selection un-made.
This is the jar's near side: everything the machine can honestly say about the coupling from inside itself, held out
plainly as a bracket and nothing more.

% [6.6.p1b]
A thought experiment --- the ancient \emph{heap paradox}, the \emph{sorites} --- holds the vague boundary open the same
way the detange does.
Take a heap of sand and take away one grain, then another; no single grain is the one whose loss turns a heap into a
non-heap, and yet a heap taken apart grain by grain is gone at the end --- the boundary is real in its effect and nowhere
in the counting.
The machine meets this at its floor: asked which candidate the value is, convergent by convergent, it finds no single
refinement that turns one answer into another, because below the count-to-three no characteristic it can decide sets them
apart.
So it does with the boundary what the sand does --- it declines to place it, detangs the forced choice, and keeps the
bracket open where no grain of counting will close it; the sharp line is on the far side, and the honest reading names its
absence rather than inventing one.

% [6.6.p2]
Derived and read, kept apart at the very edge of the number.
That the count-to-three walls the inverse coupling into an ordered bracket, its lower wall genuinely below its upper ---
that is a theorem, decided on the build; that the crossing is the surd $\sqrt{18/5}$ --- given those two measured constants --- and its continued fraction periodic
is the machine's own arithmetic; the grammar that licenses the reading is Volume 4, \emph{The Representational Grammar}.
That this bracket \emph{is} the coupling, that the coupling \emph{is} the self-energy read to a value, that $137.011$ is
the number the machine drifts toward past its floor --- these are the gloss laid over the arithmetic, marked as
gloss, and the drifted figure is carried as the machine's own artifact, never as a constant of the world.
What the machine cannot name about the coupling --- the value a laboratory measures for it, out in the world --- is the
jar's far side.
What comes next opens it, names it once, and proves the machine cannot reach it.
